% !TeX root = ../PIFB2.tex
\section{Open Problems with Decisive Next Steps}\label{sec:open_problems}

The unresolved program is narrower than deriving everything from a scalar functional. Each problem asks whether a specified operation can be justified by the actual agent dynamics, with a result or counterexample that would change the theory.

\paragraph{Predictive parent state.}
For the coupled selective-link dynamics, construct an equivariant statistic whose future record law depends only on its retained state and a controlled memory. Test the link square and generator descent, then compute the projected memory kernel. If a finite family is proposed, prove closure or quantify its error on the formation and handoff horizons. Failure of one finite family does not exclude a different family or a measure-valued parent.

\paragraph{Recursive formation.}
Extend the established fixed-temperature selection mechanism to an inherited parent dynamics without resetting its laws or supplying a preferred second partition. Separate the asymmetry in the initial preparation from the learning that amplifies it. A second episode must follow from the actual first parent, including retained attention events and hidden conditional information. The current four-cycle theorem does not settle this.

\paragraph{Emergent model variables.}
One joint bundle is already sufficient for state and model laws. The stronger problem is to derive a stable predictive coordinate and its slow evolution from an unsplit state. Accumulating sufficient statistics and exact projected memory provide examples. Local Fisher descent at an exactly represented posterior supplies no intrinsic spectral hierarchy, so a claim of slow model emergence must identify the mechanism that avoids that special case.

\paragraph{Learned selective interfaces.}
A likelihood involving low-rank maps can be unidentifiable or favor uninformative maps. Determine which records, conservation constraints, structural priors or kinetic mechanisms select a nontrivial interface without encoding the desired hierarchy in advance. Track the term caused by motion of the retention itself. Rank alone is neither a capacity nor an agency criterion.

\paragraph{Operational Lorentz geometry.}
The determinant construction gives exact $1+3$ kinematics. The missing implication is selection of a complex two-component channel, a nondegenerate collective soldering to a four-dimensional base region, and a scale convention compatible across overlaps. A physical causal cone must then be the characteristic cone of a derived propagation law. A cone of positive moment matrices alone is not a cone of allowed trajectory velocities. Einstein dynamics, physical units and coupling to matter remain additional targets.

\paragraph{Physical inertia.}
Derive position and momentum observables, a calibrated forced response and the corresponding conservation or balance law. A second derivative obtained by eliminating hidden first-order variables does not suffice. Memory and metastable lifetimes give concrete persistence mechanisms; a claim that they become inertial mass must survive frequency-dependent response tests and distinguish friction from kinetic storage.

\paragraph{Registration and the quantum question.}
Binary record kernels provide the information-gain bound. They do not derive irreversible registration, incompatibility of measurements, the Born rule or quantum amplitudes. Classical complex messages already support the Hermitian determinant construction, so complex notation is not evidence of quantum physics. A quantum extension requires its own state space, update law and empirical tests.

\paragraph{Empirical identification.}
Choose a target with accessible records, a noise model and observables sensitive to the proposed links or hierarchy. Infer fine-scale parameters with uncertainty, then predict retained-scale records without independent refitting. Specify what observation would reject the added physical identification. Internal agreement of derivations or a successful model simulation is not that comparison.

